%!Mode:: "TeX:UTF-8"
%%% Local Variables:
%%% mode: latex
%%% TeX-master: t
%%% End:

\documentclass[doctor]{thuthesis}
% \documentclass[%
%   bachelor|master|doctor|postdoctor, % mandatory option
%   xetex|pdftex|dvips|dvipdfm, % optional
%   secret,
%   openany|openright,
%   arialtoc,arialtitle]{thuthesis}

% 所有其它可能用到的包都统一放到这里了，可以根据自己的实际添加或者删除。
\usepackage{thutils}
\usepackage{bm}
\usepackage{dsfont}


\usepackage{color}
\usepackage{algorithm,algpseudocode}



\DeclareMathOperator*{\argmin}{\arg\!\min}
\DeclareMathOperator*{\argmax}{\arg\!\max}

\newcommand{\tabincell}[2]{\begin{tabular}{@{}#1@{}}#2\end{tabular}}
\newcommand{\red}{\textcolor[rgb]{1.00,0.00,0.00}}
\newcommand{\blue}{\textcolor[rgb]{0.00,0.00,1.00}}

\DeclareMathOperator{\sgn}{sgn}

\floatname{algorithm}{算法}
\makeatletter
\@addtoreset{algorithm}{chapter}% algorithm counter resets every chapter
\makeatother
\renewcommand{\thealgorithm}{\thechapter.\arabic{algorithm}}% Algorithm # is <chapter>.<algorithm>
\algnewcommand\algorithmicinput{\textbf{输入:}}
\algnewcommand\INPUT{\item[\algorithmicinput]}
\algnewcommand\algorithmicoutput{\textbf{输出:}}
\algnewcommand\OUTPUT{\item[\algorithmicoutput]}
\algnewcommand\algorithmicinit{\textbf{初始化:}}
\algnewcommand\INIT{\item[\algorithmicinit]}


% 你可以在这里修改配置文件中的定义，导言区可以使用中文。
% \def\myname{薛瑞尼}

\begin{document}

% 定义所有的eps文件在 figures 子目录下
\graphicspath{{figures/}}


%%% 封面部分
\frontmatter
%!Mode:: "TeX:UTF-8"
%%% Local Variables:
%%% mode: latex
%%% TeX-master: t
%%% End:
\secretlevel{绝密} \secretyear{2100}

\ctitle{一份简明的THUThesis模板}
% 根据自己的情况选，不用这样复杂
\makeatletter
\ifthu@bachelor\relax\else
  \ifthu@doctor
    \cdegree{工学博士}
  \else
    \ifthu@master
      \cdegree{工学硕士}
    \fi
  \fi
\fi
\makeatother


\cdepartment[电子]{电子工程系}
\cmajor{信息与通信工程}
\cauthor{匿名君}
\csupervisor{某老师教授}
% 如果没有副指导老师或者联合指导老师，把下面两行相应的删除即可。
%\cassosupervisor{陈文光教授}
%\ccosupervisor{某某某教授}
% 日期自动生成，如果你要自己写就改这个cdate
%\cdate{\CJKdigits{\the\year}年\CJKnumber{\the\month}月}

% 博士后部分
% \cfirstdiscipline{计算机科学与技术}
% \cseconddiscipline{系统结构}
% \postdoctordate{2009年7月——2011年7月}

\etitle{A Brief Template for THUThusis}
% 这块比较复杂，需要分情况讨论：
% 1. 学术型硕士
%    \edegree：必须为Master of Arts或Master of Science（注意大小写）
%              “哲学、文学、历史学、法学、教育学、艺术学门类，公共管理学科
%               填写Master of Arts，其它填写Master of Science”
%    \emajor：“获得一级学科授权的学科填写一级学科名称，其它填写二级学科名称”
% 2. 专业型硕士
%    \edegree：“填写专业学位英文名称全称”
%    \emajor：“工程硕士填写工程领域，其它专业学位不填写此项”
% 3. 学术型博士
%    \edegree：Doctor of Philosophy（注意大小写）
%    \emajor：“获得一级学科授权的学科填写一级学科名称，其它填写二级学科名称”
% 4. 专业型博士
%    \edegree：“填写专业学位英文名称全称”
%    \emajor：不填写此项
\edegree{Doctor of Philosophy}
\emajor{Information and Communication Engineering}
\eauthor{Ni Mingjun}
\esupervisor{Professor Mou Laoshi}
%\eassosupervisor{Chen Wenguang}
% 这个日期也会自动生成，你要改么？
% \edate{December, 2005}

% 定义中英文摘要和关键字
\begin{cabstract}
本文是一份简要的~THUThisis~模板。

（博士）摘要要求是 500--800 字，和关键字必须放在一页内

\end{cabstract}

\ckeywords{论文；模板；关键字；不超过五个；紫薯补丁}

\begin{eabstract}

Translation for Chinese version, do not need to be in 1 page.

\end{eabstract}

\ekeywords{thesis; template; keyword; no more than five; patch for word count}

\makecover

% 目录
\tableofcontents

% 符号对照表
%!Mode:: "TeX:UTF-8"
\begin{denotation}
\item[DSP] 数字信号处理(Digital Signal Processing)
\item[$\forall$] 任取

%\item[]

\end{denotation}



%%% 正文部分
\mainmatter
%!Mode:: "TeX:UTF-8"
%%% Local Variables:
%%% mode: latex
%%% TeX-master: t
%%% End:

\chapter{引言}

\label{cha:intro}

这篇文章介绍的主要是如何把~THUThesis~编通，以及一些注意事项。模板的进一步说明在本模板带的文件
ThuThesis\_User\_Guide\_3.0.pdf
中。

至于文章内容如何组织，请参考研究生院的说明\cite{shuoming}和师兄的论文。

至于那些随便找个模板就能编过的，咱得膜拜一下。

我是 2016 年 1 月毕业的，至少在这个时间点上该模板还有效。



%!Mode:: "TeX:UTF-8"
%%% Local Variables:
%%% mode: latex
%%% TeX-master: t
%%% End:

\chapter{编译注意事项}
\label{cha:compile}

\begin{enumerate}[1. ]
    \item 编译用 CTeX\_2.9.2.164\_Full 版本中的 pdflatex，北邮人上有，不要做更新。如果更新了编译不过就重装一下吧。如果有人搞通了 xelatex，教教我哈$\sim$；
    \item 编译正文打开 main.tex 并编译；编译书脊打开 shuji.tex 并编译；
    \item 各章节放在 data/ 下面，main.tex 中有引用，可以增删；
    \item 图放在 figures/ 下面；
    \item bib在 ref/refs.bib 文件中；
    \item 如果文章没有任何引用 \textbackslash cite 或者 \textbackslash onlinecite，编译可能会报错；
    \item 我有时会遇到编译报错一次后，再编译出很奇怪的错误的情况。如果是 main.tex 中报的错，删掉 main.out 重新编译，再不行删掉 main.aux；如果是 chapxx.tex 报的错，删掉 chapxx.aux 再重新编译即可。
\end{enumerate}







%!Mode:: "TeX:UTF-8"
%%% Local Variables:
%%% mode: latex
%%% TeX-master: t
%%% End:

\chapter{公式}
\label{cha:3}

equation 环境+\textbackslash nonumber 时会跳号（虽然编号没显示，但实际上编号+1），像这样。


\begin{equation}
{\text{It SEEMS good here.}} \nonumber
\end{equation}

\begin{equation}
{\text{But look at my equation number!!!}}
\end{equation}

\begin{equation*}
{\text{equation* seems OK.}} \nonumber
\end{equation*}

\begin{equation}
{\text{My equation number is good!!!}}
\end{equation}






遇到此问题可用align环境（貌似所有equation都可替换成align，最后做个查找替换也行）。


\begin{align}
{\text{It SEEMS good here.}} \nonumber
\end{align}

\begin{align}
{\text{And is really OK.}}
\end{align}




如果你没遇到这个问题，说明也许是我编译器的问题，或者人品问题？！



%!Mode:: "TeX:UTF-8"
%%% Local Variables:
%%% mode: latex
%%% TeX-master: t
%%% End:

\chapter{定理定义环境和证明以及字号字体}
\label{cha:4}

定义定理啥的在 ThuThesis\_User\_Guide\_3.0.pdf 中有详尽的说明。


像这样。


\begin{theorem}
我是定理。
\end{theorem}

证明我是直接抄师兄的做法了。


\noindent{\hei 证明}：

balabalabala

\noindent{\hei 证毕}。


字号字体也在 ThuThesis\_User\_Guide\_3.0.pdf 说的很好了。
刚才证明和证毕几个字已经用到了黑体。



%!Mode:: "TeX:UTF-8"
%%% Local Variables:
%%% mode: latex
%%% TeX-master: t
%%% End:

\chapter{图片}
\label{cha:5}

之前用的 subplot 不行了。这里用 subfloat。如图 \ref{fig:example}。

\begin{figure}
\centering
\subfloat[可以有子标题也可以没有]{%
  \includegraphics[width=0.4\textwidth]{1.pdf}
}
\subfloat[]{%
  \includegraphics[width=0.4\textwidth]{2.pdf}
}
\\
\subfloat[I'm 3]{%
  \includegraphics[width=0.4\textwidth]{3.pdf}
}
\hfill
\subfloat[And at last]{%
  \includegraphics[width=0.4\textwidth]{4.pdf}
}
\caption{总标题}
\label{fig:example}
\end{figure}




%!Mode:: "TeX:UTF-8"
%%% Local Variables:
%%% mode: latex
%%% TeX-master: t
%%% End:

\chapter{表格与算法}
\label{cha:6}

如果要调字号的话，直接用 \textbackslash begin\{footnotesize\} 跟 \textbackslash end\{footnotesize\}
那套东西不行。表 \ref{tb:example} 的是一个可行的方案，用不用由你。


\begin{table}
    \centering
    \caption{表格样例}
    \label{tb:example}
    \begin{tabular}{>{\wuhao}c|>{\xiaowu}c|>{\xiaowu}c|>{\xiaowu}c|>{\xiaowu}c|>{\xiaowu}c|>{\xiaowu}c|>{\xiaowu}c|>{\xiaowu}c|>{\xiaowu}c|>{\xiaowu}c}  \hline
    no1 & \multicolumn{2}{>{\xiaowu}c|}{$10^{-5}$} & \multicolumn{2}{>{\xiaowu}c|}{$2 \times 10^{-5}$} & \multicolumn{2}{>{\xiaowu}c|}{$3 \times 10^{-5}$} & \multicolumn{2}{>{\xiaowu}c|}{$4 \times 10^{-5}$} & \multicolumn{2}{>{\xiaowu}c}{$5 \times 10^{-5}$}
     \\ \hline
    something & $0$ & $0$ & $0\%$ & $0$ & $0\%$ & 0 & $0\%$ & 0 & $0\%$ & $0$
    \\ \hline
    \parbox[c][1cm][c]{2cm}{\centering 换行\\换行} & $111$ & $0$ & $111$ & $111$ & $111$ & $0\%$ & $1\%$ & $100\%$ & $2\%$ & $100\%$
    \\ \hline
    \end{tabular}
\end{table}


算法的话我用的是 algorithmicx，它可以自动处理算法编号，包括章节。main.tex 里面已经配好了，可以直接拿来用。
如果你不想用这个，请删掉 main.tex 的相关代码并加入自己的东西。算法 \ref{alg:example} 是一个示例。


\begin{algorithm}[htb!]  % 正常就h就好了，这里是为了放到一页上好看点
\caption{算法示例}
\label{alg:example}
\begin{algorithmic}[1]
\INPUT
    \Statex 输入1；
    \Statex 输入2；
\OUTPUT
    \Statex 输出1；
\INIT
    \Statex 初始化1；
    \State 步骤1； \Comment{还可以加注释哦}
    \For{$~i=2$~ {\bf to} $~N-1$~ }
        \State do something
    \EndFor
    \Repeat
    \State lalala
        \If{$a=0$}
            \State $b=0$
        \EndIf
    \Until{some condition}
    \State 步骤2。
\end{algorithmic}
\end{algorithm}


%%%%%%


\chapter{其他}
\label{chp:others}

没人规定一个文件里只能有一章，你想吧所有正文都放一个文件都没问题呀

\section{关于文字空格}

\subsection{dummy}

\subsubsection{这一级的标题不会进入目录的哦}

pdflatex 加上这个模板中英文混编的时候我是觉得特别难看了，比方说English这种，我这里是挤在一起的，数字如2000也是这样，公式如$x$也是这样，引用如\ref{chp:others}一样。

我是不喜欢这样啦，这点 Word 就做得好呀，不知 xelatex 是不是这样，或者跟模板有关啥的。
为了解决这个问题，我都加了空格。
然后，又感觉英文和中文标点之间不应该有空格呀，如果句首是英文不应该在前面有空格呀啥的，可麻烦了。

{\hei \sanhao 所以不是实在看不下去就不要加空格了吧。。。}

如果一定要加，还是比较推荐写的时候就手动加上，像我这样查找替换为主还是挺容易出错的。

我在写论文和查找替换的时候，采取了以下规则，可能稍微方便一点，但总是有漏网之鱼或者错误替换的情况出现，这时就得靠眼睛看了。

在写论文的时候注意：
\begin{enumerate}[1. ]
    \item 纯文字的话自己填上空格，如年份，中英文夹杂啥的（其实是我不会弄）。注意如果句首是英文的话，不要加空格，更不要加换行，因为一个换行也会被 tex 认为是空格；
    \item 不使用 \$\$ 写公式，统一用 align 就好了；
    \item 中文和 \$ 之间先不加空格，\textbackslash ref\{\} 前后先不加空格，之后统一替换。
\end{enumerate}

写完以后，顺序进行如表 \ref{tb:replace} 所示的替换。
{\hei 替换之前请好做好备份}。
我这里可是没做检查，不要同一个操作进行多次哦。
里面有普通的文字替换，也有正则表达式的替换，正则表达式替换可以用 sublime 啊，Notepad++ 什么的做。
其中有一个步骤是看你的公式用的是中文括号（前后不应加空格）还是英文括号（前后应该加空格），请只选做其中的一步。
注意直接在表里复制的话，粘贴出来有问题，所以请打开源代码，复制在 \verb+\verb|+ 和 \verb+|+ 之间的代码。

\begin{table}
    \centering
\caption{替换流程}
\label{tb:replace}
    \begin{tabular}{>{\xiaowu}c|>{\xiaowu}c|>{\xiaowu}c|>{\xiaowu}c|>{\xiaowu}c}  \hline
    针对人群 & 模式 & from & to  & 作用 \\\hline
    所有 & 普通 & {\$} & \$$\sim$  & 给所有行内公式加空格 \\\hline
    所有 & 正则 & \verb|(\\ref\{[^\}]*\})| & \verb*| \1 |  & 所有 ref 前后加空格 \\\hline
    所有 & 正则 & \verb*|(\\ref\{[^\}]*\}) (\([a-z]\))| & \verb*|\1\2 |  & \parbox[c][1cm][c]{4cm}{\centering 把刚才替换结果中像\\ “图 1.1~(a)”变成“图 1.1(a) ”} \\\hline
    仅中文括号 & 正则 & \verb*|（ (\\ref\{[^\}]*\}) ）| & \verb|（\1）|  & \parbox[c][1cm][c]{4cm}{\centering 删掉包含在中文括号中\\ ref 左右的空格} \\\hline
    \parbox[c][0.6cm][c]{1cm}{\centering 仅英文\\括号} & 正则 & \verb*|\( (\\ref\{[^\}]*\}) \)| & \verb*| (\1) |  & \parbox[c][1cm][c]{5cm}{\centering 把英文括号中 ref 左右\\ 的空格挪到括号外} \\\hline
    所有 & 正则 & \verb*|[ ~]([，。（；）])| & \verb|\1\2 |  & 删除标点前面的空格 \\\hline
    所有 & 正则 & \verb*|([，。（；）]\$)[ \~]| & \verb|\1\2 |  & \parbox[c][1cm][c]{5cm}{\centering 若标点后直接接公式，\\ 删除标点后面的空格}    \\ \hline
    \end{tabular}
\end{table}

\section{关于本模板附带的两个脚本文件}

本模板附带了两个 bash 脚本文件，一个是 git 一键提交（c.sh），另外一个是 git 自动 pull 并编译（auto\_pull.sh）。

因为我在 MAC 下写论文但是在 WIN 下编译，这套东西就挺有用的了。
写了一点想看结果，直接 ./c.sh （后面可以加参数表示 commit 的说明），那边 WIN 下跑的 auto\_pull.sh 就会自动拉下来编译，CTEX 自带的那个编辑器还会自动刷新 PDF。

如果不想用 winedt 做编辑的话，我觉得还是挺方便的，顺便连备份都做了。


\section{如果你正在用或者想用 VIM 写论文}

本人基本等价于标准程序猿一枚。

如果你正在用或者想在 MAC 系统下用 VIM 写论文，并且你发现 Normal 模式经常被中文输入法打扰，你可以考虑这一页的解决方案\onlinecite{macvim}。虽然不完美，但总比没有好。

WIN 下应该也有类似的解决方案，但我没找过。

才不会告诉你我的毕业论文跟这个文档都是在 VIM 下码的字呢。

%%%

\chapter{结束语}
\label{chp:ending}

就这些了，祝你顺利毕业$\sim$




%%% 其它部分
\backmatter

% 本科生要这几个索引，研究生不要。选择性留下。
\makeatletter
\ifthu@bachelor
  % 插图索引
  \listoffigures
  % 表格索引
  \listoftables
  % 公式索引
  \listofequations
\fi
\makeatother


% 参考文献
\bibliographystyle{thubib}
\bibliography{ref/refs}


% 致谢

%%% Local Variables:
%%% mode: latex
%%% TeX-master: "../main"
%%% End:

\begin{ack}
%感谢党。

感谢老师。

感谢童鞋。

感谢家人。

\end{ack}


% 附录
%\begin{appendix}
%%\chapter{外文资料原文}
%\label{cha:engorg}
%As one of the most widely used techniques in operations research, {\em
%  mathematical programming} is defined as a means of maximizing a quantity known
%as {\em objective function}, subject to a set of constraints represented by
%equations and inequalities. Some known subtopics of mathematical programming are
%linear programming, nonlinear programming, multiobjective programming, goal
%programming, dynamic programming, and multilevel programming$^{[1]}$.
%
%It is impossible to cover in a single chapter every concept of mathematical
%programming. This chapter introduces only the basic concepts and techniques of
%mathematical programming such that readers gain an understanding of them
%throughout the book$^{[2,3]}$.
%
%
%\section{Single-Objective Programming}
%The general form of single-objective programming (SOP) is written
%as follows,
%\begin{equation}\tag*{(123)} % 如果附录中的公式不想让它出现在公式索引中，那就请
%                             % 用 \tag*{xxxx}
%\left\{\begin{array}{l}
%\max \,\,f(x)\\[0.1 cm]
%\mbox{subject to:} \\ [0.1 cm]
%\qquad g_j(x)\le 0,\quad j=1,2,\cdots,p
%\end{array}\right.
%\end{equation}
%which maximizes a real-valued function $f$ of
%$x=(x_1,x_2,\cdots,x_n)$ subject to a set of constraints.
%
%\newtheorem{mpdef}{Definition}[chapter]
%\begin{mpdef}
%In SOP, we call $x$ a decision vector, and
%$x_1,x_2,\cdots,x_n$ decision variables. The function
%$f$ is called the objective function. The set
%\begin{equation}\tag*{(456)} % 这里同理，其它不再一一指定。
%S=\left\{x\in\Re^n\bigm|g_j(x)\le 0,\,j=1,2,\cdots,p\right\}
%\end{equation}
%is called the feasible set. An element $x$ in $S$ is called a
%feasible solution.
%\end{mpdef}
%
%\newtheorem{mpdefop}[mpdef]{Definition}
%\begin{mpdefop}
%A feasible solution $x^*$ is called the optimal
%solution of SOP if and only if
%\begin{equation}
%f(x^*)\ge f(x)
%\end{equation}
%for any feasible solution $x$.
%\end{mpdefop}
%
%One of the outstanding contributions to mathematical programming was known as
%the Kuhn-Tucker conditions\ref{eq:ktc}. In order to introduce them, let us give
%some definitions. An inequality constraint $g_j(x)\le 0$ is said to be active at
%a point $x^*$ if $g_j(x^*)=0$. A point $x^*$ satisfying $g_j(x^*)\le 0$ is said
%to be regular if the gradient vectors $\nabla g_j(x)$ of all active constraints
%are linearly independent.
%
%Let $x^*$ be a regular point of the constraints of SOP and assume that all the
%functions $f(x)$ and $g_j(x),j=1,2,\cdots,p$ are differentiable. If $x^*$ is a
%local optimal solution, then there exist Lagrange multipliers
%$\lambda_j,j=1,2,\cdots,p$ such that the following Kuhn-Tucker conditions hold,
%\begin{equation}
%\label{eq:ktc}
%\left\{\begin{array}{l}
%    \nabla f(x^*)-\sum\limits_{j=1}^p\lambda_j\nabla g_j(x^*)=0\\[0.3cm]
%    \lambda_jg_j(x^*)=0,\quad j=1,2,\cdots,p\\[0.2cm]
%    \lambda_j\ge 0,\quad j=1,2,\cdots,p.
%\end{array}\right.
%\end{equation}
%If all the functions $f(x)$ and $g_j(x),j=1,2,\cdots,p$ are convex and
%differentiable, and the point $x^*$ satisfies the Kuhn-Tucker conditions
%(\ref{eq:ktc}), then it has been proved that the point $x^*$ is a global optimal
%solution of SOP.
%
%\subsection{Linear Programming}
%\label{sec:lp}
%
%If the functions $f(x),g_j(x),j=1,2,\cdots,p$ are all linear, then SOP is called
%a {\em linear programming}.
%
%The feasible set of linear is always convex. A point $x$ is called an extreme
%point of convex set $S$ if $x\in S$ and $x$ cannot be expressed as a convex
%combination of two points in $S$. It has been shown that the optimal solution to
%linear programming corresponds to an extreme point of its feasible set provided
%that the feasible set $S$ is bounded. This fact is the basis of the {\em simplex
%  algorithm} which was developed by Dantzig as a very efficient method for
%solving linear programming.
%\begin{table}[ht]
%\centering
%  \centering
%  \caption*{Table~1\hskip1em This is an example for manually numbered table, which
%    would not appear in the list of tables}
%  \label{tab:badtabular2}
%  \begin{tabular}[c]{|c|m{0.8in}|c|c|c|c|c|}\hline
%    \multicolumn{2}{|c|}{Network Topology} & \# of nodes &
%    \multicolumn{3}{c|}{\# of clients} & Server \\\hline
%    GT-ITM & Waxman Transit-Stub & 600 &
%    \multirow{2}{2em}{2\%}&
%    \multirow{2}{2em}{10\%}&
%    \multirow{2}{2em}{50\%}&
%    \multirow{2}{1.2in}{Max. Connectivity}\\\cline{1-3}
%    \multicolumn{2}{|c|}{Inet-2.1} & 6000 & & & &\\\hline
%    \multirow{2}{1in}{Xue} & Rui  & Ni &\multicolumn{4}{c|}{\multirow{2}*{\thuthesis}}\\\cline{2-3}
%    & \multicolumn{2}{c|}{ABCDEF} &\multicolumn{4}{c|}{} \\\hline
%\end{tabular}
%\end{table}
%
%Roughly speaking, the simplex algorithm examines only the extreme points of the
%feasible set, rather than all feasible points. At first, the simplex algorithm
%selects an extreme point as the initial point. The successive extreme point is
%selected so as to improve the objective function value. The procedure is
%repeated until no improvement in objective function value can be made. The last
%extreme point is the optimal solution.
%
%\subsection{Nonlinear Programming}
%
%If at least one of the functions $f(x),g_j(x),j=1,2,\cdots,p$ is nonlinear, then
%SOP is called a {\em nonlinear programming}.
%
%A large number of classical optimization methods have been developed to treat
%special-structural nonlinear programming based on the mathematical theory
%concerned with analyzing the structure of problems.
%\begin{figure}[h]
%  \centering
%  \includegraphics[clip]{thu-lib-logo}
%  \caption*{Figure~1\hskip1em This is an example for manually numbered figure,
%    which would not appear in the list of figures}
%  \label{tab:badfigure2}
%\end{figure}
%
%Now we consider a nonlinear programming which is confronted solely with
%maximizing a real-valued function with domain $\Re^n$.  Whether derivatives are
%available or not, the usual strategy is first to select a point in $\Re^n$ which
%is thought to be the most likely place where the maximum exists. If there is no
%information available on which to base such a selection, a point is chosen at
%random. From this first point an attempt is made to construct a sequence of
%points, each of which yields an improved objective function value over its
%predecessor. The next point to be added to the sequence is chosen by analyzing
%the behavior of the function at the previous points. This construction continues
%until some termination criterion is met. Methods based upon this strategy are
%called {\em ascent methods}, which can be classified as {\em direct methods},
%{\em gradient methods}, and {\em Hessian methods} according to the information
%about the behavior of objective function $f$. Direct methods require only that
%the function can be evaluated at each point. Gradient methods require the
%evaluation of first derivatives of $f$. Hessian methods require the evaluation
%of second derivatives. In fact, there is no superior method for all
%problems. The efficiency of a method is very much dependent upon the objective
%function.
%
%\subsection{Integer Programming}
%
%{\em Integer programming} is a special mathematical programming in which all of
%the variables are assumed to be only integer values. When there are not only
%integer variables but also conventional continuous variables, we call it {\em
%  mixed integer programming}. If all the variables are assumed either 0 or 1,
%then the problem is termed a {\em zero-one programming}. Although integer
%programming can be solved by an {\em exhaustive enumeration} theoretically, it
%is impractical to solve realistically sized integer programming problems. The
%most successful algorithm so far found to solve integer programming is called
%the {\em branch-and-bound enumeration} developed by Balas (1965) and Dakin
%(1965). The other technique to integer programming is the {\em cutting plane
%  method} developed by Gomory (1959).
%
%\hfill\textit{Uncertain Programming\/}\quad(\textsl{BaoDing Liu, 2006.2})
%
%\section*{References}
%\noindent{\itshape NOTE: these references are only for demonstration, they are
%  not real citations in the original text.}
%
%\begin{enumerate}[{$[$}1{$]$}]
%\item Donald E. Knuth. The \TeX book. Addison-Wesley, 1984. ISBN: 0-201-13448-9
%\item Paul W. Abrahams, Karl Berry and Kathryn A. Hargreaves. \TeX\ for the
%  Impatient. Addison-Wesley, 1990. ISBN: 0-201-51375-7
%\item David Salomon. The advanced \TeX book.  New York : Springer, 1995. ISBN:0-387-94556-3
%\end{enumerate}
%
%\chapter{外文资料的调研阅读报告或书面翻译}
%\section{单目标规划}
%北冥有鱼，其名为鲲。鲲之大，不知其几千里也。化而为鸟，其名为鹏。鹏之背，不知其几
%千里也。怒而飞，其翼若垂天之云。是鸟也，海运则将徙于南冥。南冥者，天池也。
%\begin{equation}\tag*{(123)}
% p(y|\mathbf{x}) = \frac{p(\mathbf{x},y)}{p(\mathbf{x})}=
%\frac{p(\mathbf{x}|y)p(y)}{p(\mathbf{x})}
%\end{equation}
%
%吾生也有涯，而知也无涯。以有涯随无涯，殆已！已而为知者，殆而已矣！为善无近名，为
%恶无近刑，缘督以为经，可以保身，可以全生，可以养亲，可以尽年。
%
%\subsection{线性规划}
%庖丁为文惠君解牛，手之所触，肩之所倚，足之所履，膝之所倚，砉然响然，奏刀騞然，莫
%不中音，合于桑林之舞，乃中经首之会。
%\begin{table}[ht]
%\centering
%  \centering
%  \caption*{表~1\hskip1em 这是手动编号但不出现在索引中的一个表格例子}
%  \label{tab:badtabular3}
%  \begin{tabular}[c]{|c|m{0.8in}|c|c|c|c|c|}\hline
%    \multicolumn{2}{|c|}{Network Topology} & \# of nodes &
%    \multicolumn{3}{c|}{\# of clients} & Server \\\hline
%    GT-ITM & Waxman Transit-Stub & 600 &
%    \multirow{2}{2em}{2\%}&
%    \multirow{2}{2em}{10\%}&
%    \multirow{2}{2em}{50\%}&
%    \multirow{2}{1.2in}{Max. Connectivity}\\\cline{1-3}
%    \multicolumn{2}{|c|}{Inet-2.1} & 6000 & & & &\\\hline
%    \multirow{2}{1in}{Xue} & Rui  & Ni &\multicolumn{4}{c|}{\multirow{2}*{\thuthesis}}\\\cline{2-3}
%    & \multicolumn{2}{c|}{ABCDEF} &\multicolumn{4}{c|}{} \\\hline
%\end{tabular}
%\end{table}
%
%文惠君曰：“嘻，善哉！技盖至此乎？”庖丁释刀对曰：“臣之所好者道也，进乎技矣。始臣之
%解牛之时，所见无非全牛者；三年之后，未尝见全牛也；方今之时，臣以神遇而不以目视，
%官知止而神欲行。依乎天理，批大郤，导大窾，因其固然。技经肯綮之未尝，而况大坬乎！
%良庖岁更刀，割也；族庖月更刀，折也；今臣之刀十九年矣，所解数千牛矣，而刀刃若新发
%于硎。彼节者有间而刀刃者无厚，以无厚入有间，恢恢乎其于游刃必有余地矣。是以十九年
%而刀刃若新发于硎。虽然，每至于族，吾见其难为，怵然为戒，视为止，行为迟，动刀甚微，
%謋然已解，如土委地。提刀而立，为之而四顾，为之踌躇满志，善刀而藏之。”
%
%文惠君曰：“善哉！吾闻庖丁之言，得养生焉。”
%
%
%\subsection{非线性规划}
%孔子与柳下季为友，柳下季之弟名曰盗跖。盗跖从卒九千人，横行天下，侵暴诸侯。穴室枢
%户，驱人牛马，取人妇女。贪得忘亲，不顾父母兄弟，不祭先祖。所过之邑，大国守城，小
%国入保，万民苦之。孔子谓柳下季曰：“夫为人父者，必能诏其子；为人兄者，必能教其弟。
%若父不能诏其子，兄不能教其弟，则无贵父子兄弟之亲矣。今先生，世之才士也，弟为盗
%跖，为天下害，而弗能教也，丘窃为先生羞之。丘请为先生往说之。”
%\begin{figure}[h]
%  \centering
%  \includegraphics{hello}
%  \caption*{图~1\hskip1em 这是手动编号但不出现索引中的图片的例子}
%  \label{tab:badfigure3}
%\end{figure}
%
%柳下季曰：“先生言为人父者必能诏其子，为人兄者必能教其弟，若子不听父之诏，弟不受
%兄之教，虽今先生之辩，将奈之何哉？且跖之为人也，心如涌泉，意如飘风，强足以距敌，
%辩足以饰非。顺其心则喜，逆其心则怒，易辱人以言。先生必无往。”
%
%孔子不听，颜回为驭，子贡为右，往见盗跖。
%
%\subsection{整数规划}
%盗跖乃方休卒徒大山之阳，脍人肝而餔之。孔子下车而前，见谒者曰：“鲁人孔丘，闻将军
%高义，敬再拜谒者。”谒者入通。盗跖闻之大怒，目如明星，发上指冠，曰：“此夫鲁国之
%巧伪人孔丘非邪？为我告之：尔作言造语，妄称文、武，冠枝木之冠，带死牛之胁，多辞缪
%说，不耕而食，不织而衣，摇唇鼓舌，擅生是非，以迷天下之主，使天下学士不反其本，妄
%作孝弟，而侥幸于封侯富贵者也。子之罪大极重，疾走归！不然，我将以子肝益昼餔之膳。”
%
%
%\chapter{其它附录}
%前面两个附录主要是给本科生做例子。其它附录的内容可以放到这里，当然如果你愿意，可
%以把这部分也放到独立的文件中，然后将其 \verb|\input| 到主文件中。 
%\end{appendix}

% 个人简历
%!Mode:: "TeX:UTF-8"
\begin{resume}
  \resumeitem{个人简历}
  1900 年 1 月出生于 C 省 D 市。

  2000 年 9 月考入 清华大学 电子系 电子信息科学与技术专业，3000 年 7 月本科毕业并获得工学学士学位， 同时获得免试攻读清华大学电子系博士学位资格。

  3000 年 9 月进入清华大学电子系攻读博士学位至今。

  \resumeitem{发表的学术论文} % 发表的和录用的合在一起

  \begin{enumerate}[{[}1{]}]
  \item M. Ni, and L. Mou, Some Paper Written by Me,
  IEEE International Conference on Acoustics, Speech, and Signal Processing  (ICASSP),
  0000-1111, May 4-9, 2014, Florence, Italy. (EI 收录, 检索号: 2014000000000.)
  \item M. Ni, and L. Mou, Some Paper Appeared in some Journal, IEEE Transactions on Signal Processing, 11(1):1111-2222, 2015.
   (SCI 收录, 检索号: WOS:000350000000020.)
  \end{enumerate}

  \resumeitem{研究成果} % 有就写，没有就删除

  \begin{enumerate}[{[}1{]}]
  \item 匿名君，某老师. 一种不知道是用来干什么的装置和方法，CN100000000B. (中国发明专利授权公告号.)
  \end{enumerate}

\end{resume}

\end{document}
